\section{Vanishing of the main energy}
\label{sec:conclusion}

The main kernel in \eqref{eq:43} has finitely many bounded site features.
We first approximate these features by multiplicative weights and the lag
factor by a function with a fixed period.  Subdividing the position block
then reduces its energy to products of the short averages in
Lemma~\ref{lem:labels-short}.  In the vanishing argument, the regularity
parameters, $C_*$, $\eta$, $C_6$, and the coarse partition are fixed;
further approximation parameters are chosen before the limits.  We take
$x\to\infty$ along the bad subsequence fixed in Section~\ref{sec:labels},
and only then $B\to\infty$.

Choose a compact interval $\mathcal I_s=[s_0,s_1]\subset(0,\infty)$
containing the supports in $s$ of all the kernels under consideration.
It depends only on the fixed smooth functions: the support conditions
$\psi(u_1)\psi(u_2)\ne0$ and
$\phi(s/u_1)\phi(s/u_2)\ne0$ put $s$ in such a fixed interval.
For the actual site sets
\[
 S_i(n)=\{p\in\mathcal P:p\mid n+i\},
\]
write the main-kernel energy as
\begin{equation}
 \mathcal Q^{\mathrm{main}}_{B,x}
 =\frac1{Tx}\sum_{n:\,n/(Tx)\in\mathcal I_s}\frac1M
   \sum_{\substack{1\le i,k\le M\\\eta T\le |k-i|\le T}}
       \overline{F_x(n+i)}F_x(n+k)
       \mathcal M_{ik}(S_i(n),S_k(n);n/(Tx)).
 \label{eq:conclusion-main-energy}
\end{equation}
The summands are zero away from the original smooth supports.  In
particular, the choice of the enclosing interval introduces no new term.

\begin{lemma}[Approximation of the site features]
\label{lem:conclusion-features}
Fix the coarse partition $I=\bigcup_{l=1}^{m_1}I_l$ used in
\eqref{eq:43}.  For every $\zeta>0$ there are real polynomials
\[
 p_l(t)=\sum_{v=0}^{d_l}\alpha_{l,v}t^v,
 \qquad 1\le l\le m_1,
\]
independent of $B$, such that the functions
\[
 \widetilde V_{l,B}(S)
   =\sum_{v=0}^{d_l}\alpha_{l,v}
       \prod_{p\in S}\frac{1+p^{-v/B}}2
\]
satisfy, for all sufficiently large $B$,
\[
 \|V_l-\widetilde V_{l,B}\|_{L^2(S)}\le\zeta.
\]
Here the norm uses the independent site law with inclusion probabilities
$1/p$.  For $0<\zeta\le1$, the approximants can be chosen with a common
bound depending only on the fixed coarse partition.  At actual integer
sites they are the fixed finite combinations
\[
 \widetilde V_{l,B}(S_i(n))
       =\sum_{v=0}^{d_l}\alpha_{l,v}G_{B,v}(n+i),
\]
where $G_{B,v}$ is the function in \eqref{eq:labels-G} with
$P_B=\mathcal P$.
\end{lemma}

\begin{proof}
Write $x_b=(\log b)/B$ for the normalized logarithm of the fair split
product $b$, and put $q_l^0(y)=\ind_{I_l}(y)/|I_l|$ for $y\ge0$.  The endpoints of the
coarse cells lie in $[1/2,3]$, where \eqref{eq:2} bounds the marginal
split-product densities uniformly for large $B$.  Choose small endpoint
neighborhoods, still in a fixed compact subinterval of $(0,3.2)$, and
replace $q_l^0$ by a continuous compactly supported function $q_l$.
We may arrange $0\le q_l\le1/|I_l|$, with equality to $q_l^0$ outside
these neighborhoods.  By \eqref{eq:2}, their $\nu_{1/2}$ mass is at
most a constant times their total length plus $O(B^{-80})$.
Consequently the neighborhoods can be fixed so that
\[
 \limsup_{B\to\infty}
   \E_{\nu_{1/2}}|q_l^0(x_b)-q_l(x_b)|^2
\]
is as small as desired, simultaneously for the finitely many cells.

The function $t\mapsto q_l(-\log t)$ on $(0,1]$, assigned value zero
at $t=0$, is continuous on $[0,1]$.  Indeed $q_l$ is zero for all
sufficiently large arguments.  Uniform polynomial approximation on
$[0,1]$ therefore gives a real polynomial $p_l$ for which
$p_l(e^{-y})$ approximates $q_l(y)$ uniformly for every $y\ge0$.
Both this approximation and the preceding endpoint smoothing can be
chosen so that
\[
 \bigl\|q_l^0(x_b)-p_l(e^{-x_b})\bigr\|_{L^2(\nu_{1/2})}
 \le\zeta
\]
for all sufficiently large $B$.  Taking the uniform polynomial error
at most one also ensures
$\sup_{0\le t\le1}|p_l(t)|\le1/|I_l|+1$.

Conditional on a site set $S$, fair splitting gives
\[
 V_l(S)=\E_{\mathrm{split}}[q_l^0(x_b)\mid S],
 \qquad
 \E_{\mathrm{split}}[e^{-v x_b}\mid S]
      =\prod_{p\in S}\left(\frac12+\frac12p^{-v/B}\right).
\]
The second identity follows by making the fair split choices separately
for each prime.  Conditional expectation is a contraction in $L^2$,
which proves the asserted error bound and the uniform bound on the
approximants.  The last identity applies to the distinct primes dividing
$n+i$, irrespective of their multiplicities, and hence gives exactly
$G_{B,v}(n+i)$ at the integer sites.
\end{proof}

We spell out how the feature error is used against the labels.
For fixed $B$, any function of $S_i(n)$ is a function of finitely many
residues of $n+i$.  Unweighted progression counting gives, for $r=1,2$,
\begin{equation}
 \lim_{x\to\infty}\frac1{Tx}
    \sum_{n:\,n/(Tx)\in\mathcal I_s}
       |V_l(S_i(n))-\widetilde V_{l,B}(S_i(n))|^r
 =|\mathcal I_s|\,
       \E|V_l(S)-\widetilde V_{l,B}(S)|^r.
 \label{eq:conclusion-feature-counting}
\end{equation}
There are only finitely many $i$ at fixed $B$; neither a bound on the
size of the residue modulus nor uniformity with respect to growing $B$
is required for this inner limit.

The singular series in \eqref{eq:41} is nonnegative and at most
$j/\varphi(j)$.  Its ordinary mean is bounded.  Thus the lag majorant
has deterministic row bounds
\begin{equation}
 \sup_{1\le i\le M}
   \frac1T\sum_{\substack{1\le k\le M\\0<|k-i|\le T}}
      \mathfrak S(|k-i|)\ll1.
 \label{eq:conclusion-lag-rows}
\end{equation}
For completeness, the identity
\[
 \frac{j}{\varphi(j)}
   =\sum_{d\mid j}\frac{\mu_{\mathrm{Mob}}^2(d)}{\varphi(d)}
\]
and the convergence of
$\sum_d\mu_{\mathrm{Mob}}^2(d)/(d\varphi(d))$ give this bound by
summing the divisor expansion up to $T$.
Since $|F_x|\le2$, the features and their approximants are bounded,
and $k_B$, $w$ and the finitely many $h_l$ are bounded for the present
fixed parameters, replacing both features in each term of
\eqref{eq:conclusion-main-energy} costs at most $C\zeta$ in iterated
upper limit.  Indeed
\[
 |V_l(S_i)V_l(S_k)
       -\widetilde V_{l,B}(S_i)\widetilde V_{l,B}(S_k)|
 \le C_{m_1}\bigl(
      |V_l(S_i)-\widetilde V_{l,B}(S_i)|
     +|V_l(S_k)-\widetilde V_{l,B}(S_k)|\bigr),
\]
and \eqref{eq:conclusion-lag-rows} reduces the sum to the single-site
$L^1$ errors in \eqref{eq:conclusion-feature-counting}.  This uses no
independence between either feature error and the labels or the other
endpoint.

\begin{lemma}[Periodic approximation of the lag factor]
\label{lem:conclusion-periodic}
For a fixed prime cutoff $P_1\ge2$, define on all integers $j$
\[
 \mathfrak S_{P_1}(j)
 =\prod_{\substack{p\le P_1\\p\mid j}}(1-1/p)^{-1}
   \prod_{\substack{p\le P_1\\p\nmid j}}
       \left(1-\frac1{(p-1)^2}\right).
\]
This is a bounded nonnegative function of period
$D=\prod_{p\le P_1}p$.  Moreover,
\[
 \lim_{P_1\to\infty}\ 
   \sup_{\substack{T\in\N\\T\ge1}}\frac1T
       \sum_{1\le j\le T}
         |\mathfrak S(j)-\mathfrak S_{P_1}(j)|=0,
\]
with the supremum taken over positive integers $T$.
\end{lemma}

\begin{proof}
Write
\[
 a(j)=\prod_{p\mid j}(1-1/p)^{-1},\qquad
 a_{P_1}(j)=\prod_{\substack{p\le P_1\\p\mid j}}(1-1/p)^{-1}.
\]
The omitted nondividing-prime factors form a product of numbers in
$(0,1]$.  Since $P_1\ge2$, its difference from one is at most
\[
 \sum_{p>P_1}\frac1{(p-1)^2}\ll P_1^{-1},
\]
uniformly in $j$.  All the retained nondividing-prime factors are in
$[0,1]$, including the possible zero factor at $p=2$.  Consequently
\[
 |\mathfrak S(j)-\mathfrak S_{P_1}(j)|
 \le a(j)-a_{P_1}(j)+C P_1^{-1}a(j).
\]
The positive divisor expansion gives
\[
 a(j)-a_{P_1}(j)
  =\sum_{\substack{d\mid j\\
                   d\text{ has a prime divisor greater than }P_1}}
      \frac{\mu_{\mathrm{Mob}}^2(d)}{\varphi(d)}.
\]
Its mean up to $T$ is bounded by
\[
 \sum_{\substack{d\ge1\\
                   d\text{ has a prime divisor greater than }P_1}}
     \frac{\mu_{\mathrm{Mob}}^2(d)}{d\varphi(d)}.
\]
The unrestricted positive series equals
$\prod_p(1+1/(p(p-1)))<\infty$, so this tail tends to zero.
The same expansion bounds the mean of $a(j)$ by that product,
uniformly in $T$.  The required conclusion follows.
\end{proof}

\begin{proposition}[Vanishing of the main energy]
\label{prop:conclusion-main-vanishes}
With $C_*$, $\eta$, $C_6$, and the coarse partition fixed as above,
\[
 \lim_{B\to\infty}\ \limsup_{x\to\infty}
                   |\mathcal Q^{\mathrm{main}}_{B,x}|=0.
\]
\end{proposition}

\begin{proof}
Let $\gamma>0$ be arbitrary.  First choose the approximants in
Lemma~\ref{lem:conclusion-features} with sufficiently small $\zeta$
that the feature replacement has iterated upper-limit error at most
$\gamma/3$.  Their degrees and coefficients are now fixed, and the
approximants have the stated bounds depending only on the fixed coarse
partition.  Next choose a fixed $P_1$ by
Lemma~\ref{lem:conclusion-periodic}.  Replacing
$\mathfrak S(|k-i|)$ by $\mathfrak S_{P_1}(k-i)$ in the energy costs
at most $\gamma/3$ in upper limit.  To check the normalization, the
absolute error is bounded by a fixed constant times
\[
 \frac1{MT}\sum_{i=1}^M
   \sum_{\substack{1\le k\le M\\0<|k-i|\le T}}
      |\mathfrak S(|k-i|)-\mathfrak S_{P_1}(k-i)|
 \le\frac2T\sum_{j=1}^T
      |\mathfrak S(j)-\mathfrak S_{P_1}(j)|.
\]
The label and origin bounds have been absorbed in that fixed constant.

Choose a fixed $0<\delta<\eta/4$, to be made small after $P_1$ has
been fixed.  Partition $\{1,\ldots,M\}$ into a fixed number $R_0$ of
consecutive blocks $J_a$, $1\le a\le R_0$, of lengths $L_a$ differing
by at most one.  Choosing $R_0$ sufficiently large in terms of
$\delta,C_6$, for all sufficiently large $B$ we have
\[
 c_{\delta,C_6}T\le L_a\le\delta T
 \qquad(1\le a\le R_0)
\]
with a positive fixed lower constant.  In particular all these lengths
tend to infinity.

Call an ordered block pair fully allowed when every $(i,k)\in J_a\times
J_b$ satisfies $\eta T\le|k-i|\le T$.  Discard any block pair
containing both an allowed and a nonallowed pair of positions.  The
values of $|k-i|$ on a block pair vary by at most $2\delta T$.  Thus
each discarded allowed pair has lag within $2\delta T$ of $\eta T$
or $T$.  For each row there are only $O(\delta T+1)$ such positions.
The factors now present, including the periodic lag factor, are bounded
for the fixed $P_1$ and coarse partition.  With the normalization
$1/(MT)$, this discarding costs $O(\delta)+O(1/T)$, with a constant
allowed to depend on these already fixed parameters.

Choose a representative position $i_a\in J_a$ for every block, and
for every fully allowed block pair put
$\lambda_{ab}=|i_b-i_a|/T\in[\eta,1]$.  Smoothness of $w$ on the
fixed compact $s$-range and on $\eta\le\lambda\le1$ gives
\[
 \sup_{s\in\mathcal I_s}
  |w(s,|k-i|/T)-w(s,\lambda_{ab})|\ll_\eta\delta
   \qquad(i\in J_a, k\in J_b).
\]
Replacing the weight by this representative value therefore has
another $O(\delta)$ cost.  We fix $\delta$ sufficiently small that
the two block errors together have upper limit at most $\gamma/3$.
Neither the representatives nor their scaled lags have to converge as
$B$ grows; only this uniform bound is used.

It remains to show that the fully factored expression tends to zero.
For $0\le r<D$ and $1\le l\le m_1$, define
\[
 A_{a,r,l}(n)
   =\frac1{L_a}
       \sum_{\substack{i\in J_a\\i\equiv r\pmod D}}
          F_x(n+i)\widetilde V_{l,B}(S_i(n)).
\]
By the feature identity, this is a fixed finite linear combination of
\[
 \frac1{L_a}\sum_{i\in J_a}
       \ind_{i\equiv r\pmod D}F_x(n+i)G_{B,v}(n+i).
\]
The lengths $L_a$ tend to infinity and the starting indices are fixed
for each $B$.  The modulus $D$, monomials $v$, and polynomial
coefficients are all fixed before $B$ grows.  Lemma~\ref{lem:labels-short}
applies with origins
$n\in[s_0Tx,s_1Tx]$, and yields
\begin{equation}
 \lim_{B\to\infty}\ \limsup_{x\to\infty}
    \frac1{Tx}\sum_{n:\,n/(Tx)\in\mathcal I_s}
                   |A_{a,r,l}(n)|^2=0
 \label{eq:conclusion-block-mean-square}
\end{equation}
for every one of the finitely many indices.

Set $c_{r,r'}=\mathfrak S_{P_1}(r'-r)$, using the periodic extension
in Lemma~\ref{lem:conclusion-periodic}, including at the zero residue.
On a block pair, splitting each position by its residue class gives
the factored expression
\begin{equation}
 \frac{k_B}{Tx}
   \sum_{n:\,n/(Tx)\in\mathcal I_s}
   \sum_{l=1}^{m_1}h_l
   \sum_{(a,b)\ \mathrm{fully\ allowed}}
      \frac{L_aL_b}{MT}\,w(n/(Tx),\lambda_{ab})
      \sum_{r,r'=0}^{D-1}c_{r,r'}
            \overline{A_{a,r,l}(n)}A_{b,r',l}(n).
 \label{eq:conclusion-factorization}
\end{equation}
The feature polynomials are real, so the conjugation here matches
exactly the conjugation in the original energy.  The factors
$L_aL_b/(MT)$ are uniformly bounded.  The numbers of cells, blocks,
and residue classes are fixed, and $k_B$ and the representative smooth
weights are uniformly bounded.  Cauchy--Schwarz in the origin variable
and \eqref{eq:conclusion-block-mean-square} make every term in
\eqref{eq:conclusion-factorization} tend to zero in the asserted
iterated upper-limit sense.  This remains true when the bounded
coefficients depend on $n/(Tx)$ or on $B$.

The original main energy consequently has iterated upper limit at
most $\gamma$ in absolute value.  Since the approximation parameters
were fixed before the limits and $\gamma>0$ was arbitrary, the
proposition follows.
\end{proof}

\begin{proof}[Proof of Proposition~\ref{prop:mixed}]
Suppose mixed decorrelation fails for some fixed $J$ and two fixed
phase vectors.  Fix the resulting subsequence, the limiting profile $W$ and bump $\phi$
from Section~\ref{sec:labels}, and the smooth functions $\rho_0,\psi$
from Section~\ref{sec:amplification}.  Choose the regularity grid and
tolerances as required for the estimates of
Sections~\ref{sec:arithmetic} and~\ref{sec:moments}.  The amplification
then supplies the positive constant $c_5$ in \eqref{eq:15}.  This
constant is independent of every sufficiently large fixed $C_*$.

We collect the upper bounds, recording the dependencies needed to
choose parameters without a cycle.  The diagonal contribution to
$I_{2,x}/(BT)$ tends to zero in the prescribed order.  Equations
\eqref{eq:16}--\eqref{eq:19} express the remaining contribution as
the block energy, up to errors bounded in iterated upper limit by
\[
 C_0\eta+C_0/C_6.
\]
Here $C_0$ can be chosen independently of $C_*$, $\eta$ and $C_6$,
by the untruncated mean bound \eqref{eq:18}.

The root comparison \eqref{eq:21} and the sampling bound
\eqref{eq:45} allow replacement of this block kernel by
$\mathcal M$.  The resulting upper-limit error is bounded by
\[
 C_1 e^{-c_3C_*}+C_{\eta,C_6}\epsilon_1+C_2\epsilon_2.
\]
The constants $C_1,C_2$ do not depend on the later choices
$\eta,C_6$ or the coarse partition.  Indeed the truncation term in
\eqref{eq:44} has a coefficient independent of the lag range and
block length, and converting a real cut-norm bound to an energy of
labels of absolute value at most two costs only an absolute factor.
For example, splitting each label into its real and imaginary parts
gives a factor at most $16$.  Integrating over $\mathcal I_s$ costs
only its fixed length.  The finite-residue counting passage and the
uniformity in $s$ were established in Sections~\ref{sec:graph}
and~\ref{sec:sampling}.  The remaining $o(1)$ errors vanish as
$B\to\infty$ after the inner $x$-limit, with all the displayed
parameters fixed.

For any such fixed choices, Proposition~\ref{prop:conclusion-main-vanishes}
makes the main energy zero in upper limit.  Consequently
\begin{equation}
 \limsup_{B\to\infty}\ \limsup_{x\to\infty}
       \frac{I_{2,x}}{BT}
 \le C_1e^{-c_3C_*}
      +C_0\eta+C_0/C_6
      +C_{\eta,C_6}\epsilon_1+C_2\epsilon_2.
 \label{eq:conclusion-final-upper}
\end{equation}
First choose a sufficiently large fixed $C_*$ so that the first term
is less than $c_5/10$.  Next choose a sufficiently small fixed $\eta$
and a sufficiently large fixed $C_6$ so that the next two terms are
each less than $c_5/10$.  Then choose positive fixed
$\epsilon_1,\epsilon_2$ so that their terms are each less than
$c_5/10$, and choose the corresponding coarse partition in
\eqref{eq:32}.  All feature, periodic and block approximations in
Proposition~\ref{prop:conclusion-main-vanishes} are made subsequently
with these parameters fixed.  Thus \eqref{eq:conclusion-final-upper}
is strictly smaller than $c_5$, contradicting the lower bound
\eqref{eq:15}.

Finally, the bad subsequence was extracted from an arbitrary sequence
of integer scales on which
$x^{-1}\sum_{1\le n<x}\overline{g_x(n)}F_x(n+1)$ stays a fixed
positive distance from zero.  Every inner limit above uses the final
subsequence selected for the profile.  At each fixed $B$, all multipliers,
residue moduli and shifts are finite, so the argument never requires their
invariance or counting estimates while $B$ and $x$ grow simultaneously.
The contradiction rules out every original sequence of this kind.  Thus
the mixed correlation tends to zero along the
full sequence of integer scales, for every fixed pair of phase vectors.
\end{proof}
